\section{Related Work}

\subsection{Glot500 and Horizontal Scaling}

Glot500은 XLM-R이 지원하는 약 100개 언어를 head language로, 그 밖의 언어를 tail
language로 정의하고, multilingual model의 horizontal scaling을 제안한다
\citep{imani2023glot500}. 원 논문은 language-script 단위 corpus cleaning, 30,000
sentence' 이상 기준, SentencePiece unigram tokenizer 재학습, XLM-R vocabulary 확장,
continued MLM pretraining을 하나의 pipeline으로 구성한다. Glot500-m은 XLM-R vocabulary
250K에 약 151K new token을 추가해 401K vocabulary를 만들고, 600GB corpus에서 MLM으로
continued pretraining한다.

원 논문 Table 4는 Glot500-m이 tail group에서 특히 큰 개선을 보인다고 보고한다. 본
보고서에서는 그 표를 단순 baseline으로 가져오는 데 그치지 않고, 왜 본 실험이 같은
head/tail/all reporting structure를 따라야 하는지의 근거로 사용한다.

\begin{table}[t]
\centering
\caption{Glot500 original head/tail/all reference table from \citet{imani2023glot500}. Values are averages across five seeds; retrieval, tagging, classification, and roundtrip values are percentages.}
\label{tab:glot500-reference}
\resizebox{\textwidth}{!}{%
\begin{tabular}{lrrrrrrrrr}
\toprule
 & \multicolumn{3}{c}{tail} & \multicolumn{3}{c}{head} & \multicolumn{3}{c}{all} \\
\cmidrule(lr){2-4}\cmidrule(lr){5-7}\cmidrule(lr){8-10}
Metric & XLM-R-B & XLM-R-L & Glot500-m & XLM-R-B & XLM-R-L & Glot500-m & XLM-R-B & XLM-R-L & Glot500-m \\
\midrule
Pseudoperplexity & 304.2 & 168.6 & 12.2 & 12.5 & 8.4 & 11.8 & 247.8 & 136.4 & 11.6 \\
Sentence Retrieval Tatoeba & 32.6 & 33.6 & 59.8 & 66.2 & 71.1 & 75.0 & 56.6 & 60.4 & 70.7 \\
Sentence Retrieval Bible & 7.4 & 7.1 & 43.2 & 54.2 & 58.3 & 59.0 & 19.3 & 20.1 & 47.3 \\
Text Classification & 13.7 & 13.9 & 46.6 & 51.3 & 60.5 & 54.7 & 23.3 & 25.8 & 48.7 \\
NER & 47.5 & 51.8 & 60.7 & 61.8 & 66.0 & 63.9 & 55.3 & 59.5 & 62.4 \\
POS & 41.7 & 43.5 & 62.3 & 76.4 & 78.4 & 76.0 & 65.8 & 67.7 & 71.8 \\
Roundtrip Alignment & 2.6 & 3.1 & 4.5 & 3.4 & 4.1 & 5.5 & 2.8 & 3.3 & 4.7 \\
\bottomrule
\end{tabular}}
\end{table}


\subsection{SentencePiece Unigram and Vocabulary Extension}

SentencePiece는 library 이름이고, 그 안에서 BPE 또는 unigram language model을 사용할 수
있다 \citep{kudo2018sentencepiece}. XLM-R과 Glot500은 SentencePiece unigram 계열을 사용한다.
Unigram LM은 가능한 subword 후보 집합을 크게 시작한 뒤, EM으로 token probability를
추정하고 likelihood contribution이 낮은 token을 pruning하는 top-down 방식이다. 이 선택은
BPE처럼 greedy merge만 누적하는 방식과 달리 여러 segmentation 후보를 확률적으로 다룬다.

본 프로젝트의 tokenizer는 이 Glot500 흐름을 따른다. \codepath{tokenization/run.py}는
\codepath{--model_type=unigram}, \codepath{--byte_fallback=true},
\codepath{--character_coverage=1.0}으로 새 SentencePiece model을 학습하고, 기존 XLM-R
SentencePiece protobuf에 없는 piece만 append한다.

\subsection{Embedding Initialization for New Tokens}

Vocabulary expansion 연구에서는 새 token을 추가한 뒤 embedding을 어떻게 초기화할지가
중요하다. Yamaguchi et al.은 30K sentences, 약 0.01GB target language text만 있는
low-resource setting에서 vocabulary expansion과 target parameter initialization을 비교하고,
high-resource에서 자주 쓰이는 random initialization이 항상 최선이 아니라고 보고한다
\citep{yamaguchi2025vocab}. Mean initialization은 기존 embedding centroid를 사용하고,
FVT 계열은 새 token surface를 source tokenizer로 다시 분해한 뒤 source subtoken embedding을
평균한다 \citep{hewitt2021newembeddings,vocabextensiontutorial2026}.

본 연구는 decoder LLM이 아니라 XLM-R encoder 기반 MLM setting이지만, 질문은 같다. 새
token row를 완전히 무작위로 시작할지, 기존 XLM-R tokenizer와 embedding space가 이미 가진
표면형 정보를 재사용할지 비교한다.
